\chapter{VALIDASI DAN APLIKASI PROGRAM OPTIMASI}
\label{bab:validasi}

\section{Validasi Program Optimasi}

Dalam melakukan validasi program optimasi, digunakan hasil perhitungan
beton bertulang pada struktur portal yang dikerjakan oleh Ir. Harsoyo
dalam bukunya berjudul Gedung Bertingkat II (Harsoyo, 1976), seri bina
bangunan konstruksi beton bertulang. Geometri bangunan yang dimodelkan
sebagai struktur portal bidang digambarkan dalam gambar 4-1.

\gambarhilang{Geometri struktur portal bidang menurut Harsoyo (1976) yang
dipakai sebagai kasus validasi.}{Geometri Struktur Portal Bidang}

Langkah pertama dalam analisa struktur setelah memodelkan struktur adalah
melakukan penomoran titik dan batang, dalam kasus ini terdapat 16 titik
kumpul dan 21 batang yang terdiri dari 12 kolom dan 9 balok. Penomoran
titik dan batang dapat dilihat dalam gambar 4-2.

\gambarhilang{Penomoran titik kumpul (1--16) dan batang (1--21, terdiri
dari 12 kolom dan 9 balok) pada struktur portal bidang validasi.}{Penomoran Titik Kumpul dan Batang Portal Bidang}

Beban yang bekerja pada struktur portal bidang tersebut berupa beban
merata pada balok dan beban akibat tekanan angin pada titik 5, 9, 13 ke
arah $-x$, besarnya beban merata pada batang diperlihatkan dalam tabel
4-1. Sedangkan besarnya beban akibat tekanan angin diperlihatkan dalam
tabel 4-2.

\begin{table}[htbp]
\centering
\caption{Beban Merata Pada Struktur Portal Bidang}
\label{tab:4-1}
\begin{tabular}{cccc}
\toprule
Nomor Batang & Beban Merata Total (t/m) & Berat Sendiri (t/m) & Beban Luar (t/m) \\
\midrule
13 & 4,100 & 0,360 & 3,740 \\
14 & 3,491 & 0,360 & 3,131 \\
15 & 4,100 & 0,360 & 3,740 \\
16 & 4,100 & 0,360 & 3,740 \\
17 & 3,491 & 0,360 & 3,131 \\
18 & 4,100 & 0,360 & 3,740 \\
19 & 1,762 & 0,300 & 1,462 \\
20 & 1,447 & 0,300 & 1,147 \\
21 & 1,762 & 0,300 & 1,462 \\
\bottomrule
\end{tabular}
\end{table}

\begin{table}[htbp]
\centering
\caption{Beban Tekanan Angin}
\label{tab:4-2}
\begin{tabular}{cc}
\toprule
Nomor Titik & Beban Tekanan Angin (kg) \\
\midrule
5  & 400 \\
9  & 800 \\
13 & 800 \\
\bottomrule
\end{tabular}
\end{table}

Dalam tugas akhir ini struktur portal bidang pada gambar 4-1 dioptimasi
dengan program optimasi beton bertulang pada struktur portal ruang, dan
hasil optimasinya dibandingkan dengan hasil yang telah dikerjakan oleh
Ir. Harsoyo D. (1976), semua data masukan dapat dilihat dalam lampiran
24, dengan menggunakan mutu beton 17,5 MPa dan mutu baja tulangan 240
MPa, tekuk pada kolom dianalisa. Perbandingan hasil akhir variabel desain
antara hasil yang dikerjakan oleh Ir. Harsoyo dan hasil optimasi program
optimasi beton bertulang pada struktur portal ruang disajikan pada tabel
4-3.

\begin{landscape}
\begingroup
\footnotesize
\begin{longtable}{clcccc}
\caption{Perbandingan Volume Beton dan Berat Besi}
\label{tab:4-3}\\
\toprule
\textbf{Batang} & \textbf{Variabel Desain} & \multicolumn{2}{c}{\textbf{Harsoyo (1976)}} & \multicolumn{2}{c}{\textbf{Naftali (1999)}} \\
\midrule
\endfirsthead
\multicolumn{6}{c}{\textit{(lanjutan Tabel 4-3)}}\\
\toprule
\textbf{Batang} & \textbf{Variabel Desain} & \multicolumn{2}{c}{\textbf{Harsoyo (1976)}} & \multicolumn{2}{c}{\textbf{Naftali (1999)}} \\
\midrule
\endhead
\bottomrule
\endfoot
\bottomrule
\endlastfoot
1 & Lebar / Tinggi & \multicolumn{2}{c}{0,30 m / 0,75 m} & \multicolumn{2}{c}{0,41 m / 0,41 m} \\
  & Volume beton & \multicolumn{2}{c}{0,9 m$^3$} & \multicolumn{2}{c}{0,6724 m$^3$} \\
  & Tulangan utama & \multicolumn{2}{c}{$10\phi25/4\phi16$ (187,929 kg)} & \multicolumn{2}{c}{$8\phi25$ (136,8158 kg)} \\
  & Tulangan sengkang & \multicolumn{2}{c}{$\phi8$--150} & \multicolumn{2}{c}{$\phi12$--300} \\
2 & Lebar / Tinggi & \multicolumn{2}{c}{0,30 m / 0,75 m} & \multicolumn{2}{c}{0,41 m / 0,41 m} \\
  & Volume beton & \multicolumn{2}{c}{0,9 m$^3$} & \multicolumn{2}{c}{0,6724 m$^3$} \\
  & Tulangan utama & \multicolumn{2}{c}{$10\phi25/4\phi16$ (187,929 kg)} & \multicolumn{2}{c}{$8\phi22$ (109,0117 kg)} \\
  & Tulangan sengkang & \multicolumn{2}{c}{$\phi8$--150} & \multicolumn{2}{c}{$\phi12$--300} \\
3 & Lebar / Tinggi & \multicolumn{2}{c}{0,30 m / 0,75 m} & \multicolumn{2}{c}{0,40 m / 0,40 m} \\
  & Volume beton & \multicolumn{2}{c}{0,9 m$^3$} & \multicolumn{2}{c}{0,64 m$^3$} \\
  & Tulangan utama & \multicolumn{2}{c}{$10\phi25/4\phi16$ (187,929 kg)} & \multicolumn{2}{c}{$8\phi22$ (101,2778 kg)} \\
  & Tulangan sengkang & \multicolumn{2}{c}{$\phi8$--150} & \multicolumn{2}{c}{$\phi8$--300} \\
4 & Lebar / Tinggi & \multicolumn{2}{c}{0,30 m / 0,75 m} & \multicolumn{2}{c}{0,43 m / 0,43 m} \\
  & Volume beton & \multicolumn{2}{c}{0,9 m$^3$} & \multicolumn{2}{c}{0,7396 m$^3$} \\
  & Tulangan utama & \multicolumn{2}{c}{$10\phi25/4\phi16$ (187,929 kg)} & \multicolumn{2}{c}{$8\phi19$ (85,6326 kg)} \\
  & Tulangan sengkang & \multicolumn{2}{c}{$\phi8$--150} & \multicolumn{2}{c}{$\phi12$--300} \\
5 & Lebar / Tinggi & \multicolumn{2}{c}{0,30 m / 0,60 m} & \multicolumn{2}{c}{0,44 m / 0,44 m} \\
  & Volume beton & \multicolumn{2}{c}{0,72 m$^3$} & \multicolumn{2}{c}{0,7744 m$^3$} \\
  & Tulangan utama & \multicolumn{2}{c}{$10\phi22/4\phi16$ (137,3734 kg)} & \multicolumn{2}{c}{$8\phi20$ (93,7609 kg)} \\
  & Tulangan sengkang & \multicolumn{2}{c}{$\phi8$--180} & \multicolumn{2}{c}{$\phi12$--300} \\
6 & Lebar / Tinggi & \multicolumn{2}{c}{0,30 m / 0,60 m} & \multicolumn{2}{c}{0,44 m / 0,44 m} \\
  & Volume beton & \multicolumn{2}{c}{0,72 m$^3$} & \multicolumn{2}{c}{0,7744 m$^3$} \\
  & Tulangan utama & \multicolumn{2}{c}{$8\phi22/4\phi16$ (113,0698 kg)} & \multicolumn{2}{c}{$8\phi19$ (77,8014 kg)} \\
  & Tulangan sengkang & \multicolumn{2}{c}{$\phi8$--180} & \multicolumn{2}{c}{$\phi8$--300} \\
7 & Lebar / Tinggi & \multicolumn{2}{c}{0,30 m / 0,60 m} & \multicolumn{2}{c}{0,42 m / 0,42 m} \\
  & Volume beton & \multicolumn{2}{c}{0,72 m$^3$} & \multicolumn{2}{c}{0,7056 m$^3$} \\
  & Tulangan utama & \multicolumn{2}{c}{$8\phi22/4\phi16$ (113,0698 kg)} & \multicolumn{2}{c}{$8\phi19$ (85,1948 kg)} \\
  & Tulangan sengkang & \multicolumn{2}{c}{$\phi8$--180} & \multicolumn{2}{c}{$\phi12$--300} \\
8 & Lebar / Tinggi & \multicolumn{2}{c}{0,30 m / 0,60 m} & \multicolumn{2}{c}{0,42 m / 0,42 m} \\
  & Volume beton & \multicolumn{2}{c}{0,72 m$^3$} & \multicolumn{2}{c}{0,7056 m$^3$} \\
  & Tulangan utama & \multicolumn{2}{c}{$10\phi22/4\phi16$ (137,3734 kg)} & \multicolumn{2}{c}{$8\phi22$ (109,4494 kg)} \\
  & Tulangan sengkang & \multicolumn{2}{c}{$\phi8$--180} & \multicolumn{2}{c}{$\phi12$--300} \\
9 & Lebar / Tinggi & \multicolumn{2}{c}{0,30 m / 0,60 m} & \multicolumn{2}{c}{0,41 m / 0,41 m} \\
  & Volume beton & \multicolumn{2}{c}{0,72 m$^3$} & \multicolumn{2}{c}{0,6724 m$^3$} \\
  & Tulangan utama & \multicolumn{2}{c}{$6\phi22/4\phi16$ (88,73483 kg)} & \multicolumn{2}{c}{$8\phi20$ (92,4476 kg)} \\
  & Tulangan sengkang & \multicolumn{2}{c}{$\phi8$--180} & \multicolumn{2}{c}{$\phi12$--300} \\
10 & Lebar / Tinggi & \multicolumn{2}{c}{0,30 m / 0,60 m} & \multicolumn{2}{c}{0,45 m / 0,45 m} \\
   & Volume beton & \multicolumn{2}{c}{0,72 m$^3$} & \multicolumn{2}{c}{0,81 m$^3$} \\
   & Tulangan utama & \multicolumn{2}{c}{$6\phi22/4\phi16$ (88,73483 kg)} & \multicolumn{2}{c}{$8\phi29$ (86,5081 kg)} \\
   & Tulangan sengkang & \multicolumn{2}{c}{$\phi8$--180} & \multicolumn{2}{c}{$\phi12$--300} \\
11 & Lebar / Tinggi & \multicolumn{2}{c}{0,30 m / 0,60 m} & \multicolumn{2}{c}{0,46 m / 0,46 m} \\
   & Volume beton & \multicolumn{2}{c}{0,72 m$^3$} & \multicolumn{2}{c}{0,8464 m$^3$} \\
   & Tulangan utama & \multicolumn{2}{c}{$6\phi22/4\phi16$ (88,73483 kg)} & \multicolumn{2}{c}{$8\phi28$ (161,6033 kg)} \\
   & Tulangan sengkang & \multicolumn{2}{c}{$\phi8$--180} & \multicolumn{2}{c}{$\phi8$--300} \\
12 & Lebar / Tinggi & \multicolumn{2}{c}{0,30 m / 0,60 m} & \multicolumn{2}{c}{0,44 m / 0,44 m} \\
   & Volume beton & \multicolumn{2}{c}{0,72 m$^3$} & \multicolumn{2}{c}{0,7744 m$^3$} \\
   & Tulangan utama & \multicolumn{2}{c}{$6\phi22/4\phi16$ (88,73483 kg)} & \multicolumn{2}{c}{$4\phi25$ (76,5066 kg)} \\
   & Tulangan sengkang & \multicolumn{2}{c}{$\phi8$--180} & \multicolumn{2}{c}{$\phi12$--300} \\
13 & Lebar / Tinggi & \multicolumn{2}{c}{0,25 m / 0,65 m} & \multicolumn{2}{c}{0,33 m / 0,57 m} \\
   & Volume beton & \multicolumn{2}{c}{0,975 m$^3$} & \multicolumn{2}{c}{1,1286 m$^3$} \\
   & Tulangan utama & \multicolumn{2}{c}{$3\phi25/3\phi25/2\phi25/4\phi25/3\phi25/3\phi25$ (228,126 kg)} & \multicolumn{2}{c}{$2\phi25/3\phi28/3\phi25/2\phi22/2\phi25/3\phi28$ (212,9059 kg)} \\
   & Tulangan sengkang & \multicolumn{2}{c}{$\phi8$--100} & \multicolumn{2}{c}{$\phi12$--130} \\
14 & Lebar / Tinggi & \multicolumn{2}{c}{0,25 m / 0,65 m} & \multicolumn{2}{c}{0,33 m / 0,52 m} \\
   & Volume beton & \multicolumn{2}{c}{0,65 m$^3$} & \multicolumn{2}{c}{0,6864 m$^3$} \\
   & Tulangan utama & \multicolumn{2}{c}{$3\phi25/3\phi25/3\phi25/2\phi25/3\phi25/3\phi25$ (143,9139 kg)} & \multicolumn{2}{c}{$2\phi22/3\phi22/3\phi25/3\phi20/2\phi22/3\phi22$ (111,5333 kg)} \\
   & Tulangan sengkang & \multicolumn{2}{c}{$\phi8$--100} & \multicolumn{2}{c}{$\phi12$--150} \\
15 & Lebar / Tinggi & \multicolumn{2}{c}{0,25 m / 0,65 m} & \multicolumn{2}{c}{0,34 m / 0,55 m} \\
   & Volume beton & \multicolumn{2}{c}{0,975 m$^3$} & \multicolumn{2}{c}{1,122 m$^3$} \\
   & Tulangan utama & \multicolumn{2}{c}{$3\phi25/3\phi25/2\phi25/4\phi25/3\phi25/3\phi25$ (228,126 kg)} & \multicolumn{2}{c}{$2\phi22/3\phi28/3\phi22/2\phi22/2\phi22/3\phi28$ (190,1989 kg)} \\
   & Tulangan sengkang & \multicolumn{2}{c}{$\phi8$--100} & \multicolumn{2}{c}{$\phi12$--140} \\
16 & Lebar / Tinggi & \multicolumn{2}{c}{0,25 m / 0,65 m} & \multicolumn{2}{c}{0,31 m / 0,55 m} \\
   & Volume beton & \multicolumn{2}{c}{0,975 m$^3$} & \multicolumn{2}{c}{1,023 m$^3$} \\
   & Tulangan utama & \multicolumn{2}{c}{$3\phi25/3\phi25/2\phi25/4\phi25/3\phi25/3\phi25$ (228,126 kg)} & \multicolumn{2}{c}{$2\phi28/3\phi28/3\phi25/2\phi28/2\phi28/3\phi28$ (225,0688 kg)} \\
   & Tulangan sengkang & \multicolumn{2}{c}{$\phi8$--100} & \multicolumn{2}{c}{$\phi12$--130} \\
17 & Lebar / Tinggi & \multicolumn{2}{c}{0,25 m / 0,65 m} & \multicolumn{2}{c}{0,33 m / 0,52 m} \\
   & Volume beton & \multicolumn{2}{c}{0,65 m$^3$} & \multicolumn{2}{c}{0,6864 m$^3$} \\
   & Tulangan utama & \multicolumn{2}{c}{$3\phi25/3\phi25/3\phi25/2\phi25/3\phi25/3\phi25$ (143,9139 kg)} & \multicolumn{2}{c}{$2\phi25/3\phi25/3\phi20/2\phi25/2\phi25/3\phi25$ (111,4915 kg)} \\
   & Tulangan sengkang & \multicolumn{2}{c}{$\phi8$--100} & \multicolumn{2}{c}{$\phi12$--170} \\
18 & Lebar / Tinggi & \multicolumn{2}{c}{0,25 m / 0,65 m} & \multicolumn{2}{c}{0,34 m / 0,54 m} \\
   & Volume beton & \multicolumn{2}{c}{0,975 m$^3$} & \multicolumn{2}{c}{1,1016 m$^3$} \\
   & Tulangan utama & \multicolumn{2}{c}{$3\phi25/3\phi25/2\phi25/4\phi25/3\phi25/3\phi25$ (228,126 kg)} & \multicolumn{2}{c}{$2\phi20/3\phi28/3\phi22/2\phi25/2\phi20/3\phi28$ (191,4206 kg)} \\
   & Tulangan sengkang & \multicolumn{2}{c}{$\phi8$--100} & \multicolumn{2}{c}{$\phi12$--140} \\
19 & Lebar / Tinggi & \multicolumn{2}{c}{0,25 m / 0,50 m} & \multicolumn{2}{c}{0,33 m / 0,49 m} \\
   & Volume beton & \multicolumn{2}{c}{0,75 m$^3$} & \multicolumn{2}{c}{0,9702 m$^3$} \\
   & Tulangan utama & \multicolumn{2}{c}{$3\phi19/4\phi19/2\phi19/4\phi19/3\phi19/4\phi19$ (179,4247 kg)} & \multicolumn{2}{c}{$2\phi25/3\phi28/3\phi22/2\phi19/2\phi25/3\phi28$ (172,2049 kg)} \\
   & Tulangan sengkang & \multicolumn{2}{c}{$\phi8$--100} & \multicolumn{2}{c}{$\phi12$--160} \\
20 & Lebar / Tinggi & \multicolumn{2}{c}{0,25 m / 0,50 m} & \multicolumn{2}{c}{0,33 m / 0,49 m} \\
   & Volume beton & \multicolumn{2}{c}{0,5 m$^3$} & \multicolumn{2}{c}{0,6468 m$^3$} \\
   & Tulangan utama & \multicolumn{2}{c}{$3\phi19/2\phi19/3\phi19/2\phi19/3\phi19/2\phi19$ (91,3503 kg)} & \multicolumn{2}{c}{$2\phi22/3\phi22/3\phi25/2\phi28/2\phi22/3\phi22$ (110,299 kg)} \\
   & Tulangan sengkang & \multicolumn{2}{c}{$\phi8$--100} & \multicolumn{2}{c}{$\phi12$--180} \\
21 & Lebar / Tinggi & \multicolumn{2}{c}{0,25 m / 0,50 m} & \multicolumn{2}{c}{0,33 m / 0,51 m} \\
   & Volume beton & \multicolumn{2}{c}{0,75 m$^3$} & \multicolumn{2}{c}{1,0098 m$^3$} \\
   & Tulangan utama & \multicolumn{2}{c}{$3\phi19/4\phi19/2\phi19/4\phi19/3\phi19/4\phi19$ (179,4247 kg)} & \multicolumn{2}{c}{$2\phi22/3\phi20/3\phi25/2\phi19/2\phi22/3\phi20$ (151,9952 kg)} \\
   & Tulangan sengkang & \multicolumn{2}{c}{$\phi8$--100} & \multicolumn{2}{c}{$\phi12$--170} \\
\midrule
\multicolumn{2}{l}{\textbf{Total Volume Beton}} & \multicolumn{2}{c}{\textbf{16,56 m$^3$}} & \multicolumn{2}{c}{\textbf{17,1624 m$^3$}} \\
\multicolumn{2}{l}{\textbf{Total Berat Besi}} & \multicolumn{2}{c}{\textbf{3113,5654 kg}} & \multicolumn{2}{c}{\textbf{2693,1281 kg}} \\
\end{longtable}
\endgroup
\end{landscape}

Selanjutnya dengan menetapkan harga satuan beton Rp.\ 250.000,00/m$^3$ dan
harga satuan besi Rp.\ 5.000,00/kg dapat dihitung harga struktur secara
keseluruhan seperti yang diperlihatkan dalam tabel 4-4, dalam hal ini
harga beton dan harga besi tersebut sudah termasuk biaya pemasangan
seperti upah pekerja, kawat pengikat dan papan acuan (\textit{form
work}).

Dari tabel 4-4 ternyata harga struktur secara keseluruhan dapat dikurangi
sebesar Rp.\ 1.951.587,00 atau sebesar 9,9~\% dengan melakukan optimasi.
Proses optimasi menggunakan 168 variabel desain dan 171 struktur desain
sebagai titik coba. Waktu proses optimasi yang diperlukan adalah 117
detik menggunakan komputer dengan prosesor Pentium Celeron 333 MHz dan
RAM 64 MB.

\begin{table}[htbp]
\centering
\caption{Perbandingan Harga Struktur}
\label{tab:4-4}
\begin{tabular}{lccccc}
\toprule
Material & Harga Satuan (Rp.) & \multicolumn{2}{c}{Harsoyo (1976)} & \multicolumn{2}{c}{Naftali (1999)} \\
\cmidrule(lr){3-4}\cmidrule(lr){5-6}
 & & Pemakaian & Harga (Rp.) & Pemakaian & Harga (Rp.) \\
\midrule
Beton & 250.000,00/m$^3$ & 16,56 m$^3$ & 4.140.000 & 17,1624 m$^3$ & 4.290.600 \\
Baja  & 5.000,00/kg      & 3113,5654 kg & 15.567.827 & 2693,1281 kg & 13.465.640 \\
\midrule
\textbf{Total} & & & \textbf{19.707.827} & & \textbf{17.756.240} \\
\bottomrule
\end{tabular}
\end{table}

\section{Aplikasi Program Optimasi}

Suatu struktur portal ruang seperti pada gambar 4-3 terbuat dari beton
bertulang memiliki data struktur sebagai berikut:
\begin{enumerate}[label=\arabic*.]
\item Jumlah batang 8 buah, terdiri dari 4 buah balok dan ditopang oleh 4
      buah kolom.
\item Tumpuan jepit pada titik kumpul 1, 2, 3, 4.
\item Kuat desak karakteristik beton = 30 MPa.
\item Kuat tarik baja tulangan memanjang = 400 MPa.
\item Kuat tarik baja tulangan sengkang = 240 MPa.
\item Selimut beton pada balok = 50 mm.
\item Selimut beton pada kolom = 50 mm.
\item Beban merata sebesar 35 kN/m pada semua balok.
\item Variabel desain yang digunakan terdapat pada lampiran.
\item Faktor penalti diambil $1\times10^{10}$.
\item Jumlah struktur desain sebagai titik coba diambil sama dengan
      jumlah variabel desain $+\,3$, jumlah variabel desain $\times\,2$
      dan jumlah variabel desain $\times\,3$.
\end{enumerate}

\gambarhilang{Geometri struktur portal ruang untuk kasus aplikasi
program.}{Geometri Struktur Portal Ruang}

Penomoran titik kumpul (\textit{joint}) dan penomoran elemen batang
diperlihatkan pada gambar 4-4, pada gambar 4-4 juga diperlihatkan
pembebanan pada struktur.

\gambarhilang{Penomoran titik kumpul dan elemen batang, serta pembebanan,
pada struktur portal ruang kasus aplikasi.}{Penomoran Titik dan Batang Portal Ruang}

Program dijalankan sebanyak 15 kali dengan 3 variasi jumlah struktur
desain, yaitu jumlah variabel desain $+\,3$, jumlah variabel desain
$\times\,2$, jumlah variabel desain $\times\,3$, masing-masing sebanyak
5 kali, dengan menggunakan komputer dengan prosesor Pentium Celeron
333 MHz dan RAM 64 MB, hasil optimasi disajikan dalam bentuk grafik pada
gambar 4-5.

\gambarhilang{Sumbu tegak: nilai fitness; sumbu datar: jumlah struktur
desain sebagai titik coba (tiga variasi: $JVD+3$, $JVD\times2$,
$JVD\times3$, masing-masing 5 pengulangan).}{Grafik Perbandingan Nilai Fitness}

Grafik pada gambar 4-5 menunjukkan bahwa semakin banyak struktur desain
yang digunakan, nilai fitness yang dicapai semakin tinggi, hal ini
berarti untuk mendapatkan hasil yang makin optimal dibutuhkan jumlah
struktur desain yang makin banyak, akan tetapi dengan bertambahnya jumlah
struktur desain, waktu yang dibutuhkan untuk proses optimasi juga
semakin banyak. Grafik perbandingan waktu proses optimasi yang
dibutuhkan oleh masing-masing penambahan jumlah variabel desain
ditunjukkan dalam grafik yang terdapat pada gambar 4-6.

\gambarhilang{Sumbu tegak: waktu proses optimasi (detik); sumbu datar:
jumlah struktur desain sebagai titik coba.}{Grafik Perbandingan Lama Proses Optimasi}

Hasil terbaik untuk kasus portal ruang pada gambar 4-3 dengan nilai
fitness $1174{,}59$ memerlukan waktu proses optimasi sebesar 74 detik.
Fungsi sasarannya yaitu harga struktur, dengan mengambil harga beton
Rp.\ 250.000,00/m$^3$ dan harga besi Rp.\ 5.000,00/kg adalah sebesar
Rp.\ 8.535.430,00, hasilnya ditunjukkan pada tabel 4-5, sedangkan
perhitungan rincian harga terdapat pada tabel 4-6.

\begingroup
\footnotesize
\begin{longtable}{clcc}
\caption{Dimensi Struktur Portal Ruang Hasil Optimasi}
\label{tab:4-5}\\
\toprule
\textbf{Batang} & \textbf{Variabel Desain} & \multicolumn{2}{c}{\textbf{Penggunaan Material}} \\
\midrule
\endfirsthead
\multicolumn{4}{c}{\textit{(lanjutan Tabel 4-5)}}\\
\toprule
\textbf{Batang} & \textbf{Variabel Desain} & \multicolumn{2}{c}{\textbf{Penggunaan Material}} \\
\midrule
\endhead
\bottomrule
\endfoot
\bottomrule
\endlastfoot
1 & Lebar / Tinggi & \multicolumn{2}{c}{0,45 m / 0,45 m (1,0125 m$^3$)} \\
  & Tulangan utama & \multicolumn{2}{c}{$8\phi19$ (102,4988 kg)} \\
  & Tulangan sengkang & \multicolumn{2}{c}{$\phi10$--300} \\
2 & Lebar / Tinggi & \multicolumn{2}{c}{0,50 m / 0,50 m (1,25 m$^3$)} \\
  & Tulangan utama & \multicolumn{2}{c}{$8\phi22$ (134,7477 kg)} \\
  & Tulangan sengkang & \multicolumn{2}{c}{$\phi10$--300} \\
3 & Lebar / Tinggi & \multicolumn{2}{c}{0,50 m / 0,50 m (1,25 m$^3$)} \\
  & Tulangan utama & \multicolumn{2}{c}{$8\phi28$ (208,6947 kg)} \\
  & Tulangan sengkang & \multicolumn{2}{c}{$\phi10$--300} \\
4 & Lebar / Tinggi & \multicolumn{2}{c}{0,45 m / 0,45 m (1,0125 m$^3$)} \\
  & Tulangan utama & \multicolumn{2}{c}{$8\phi22$ (127,9512 kg)} \\
  & Tulangan sengkang & \multicolumn{2}{c}{$\phi8$--300} \\
5 & Lebar / Tinggi & \multicolumn{2}{c}{0,25 m / 0,60 m (0,9 m$^3$)} \\
  & Tulangan utama & \multicolumn{2}{c}{$4\phi25/3\phi19/4\phi19/3\phi19/4\phi25/3\phi19$ (164,686 kg)} \\
  & Tulangan sengkang & \multicolumn{2}{c}{$\phi10$--160} \\
6 & Lebar / Tinggi & \multicolumn{2}{c}{0,30 m / 0,60 m (1,08 m$^3$)} \\
  & Tulangan utama & \multicolumn{2}{c}{$2\phi25/4\phi22/4\phi22/3\phi29/2\phi25/4\phi22$ (195,3126 kg)} \\
  & Tulangan sengkang & \multicolumn{2}{c}{$\phi10$--160} \\
7 & Lebar / Tinggi & \multicolumn{2}{c}{0,35 m / 0,50 m (1,05 m$^3$)} \\
  & Tulangan utama & \multicolumn{2}{c}{$2\phi19/3\phi25/4\phi28/2\phi22/2\phi19/3\phi25$ (171,1133 kg)} \\
  & Tulangan sengkang & \multicolumn{2}{c}{$\phi10$--160} \\
8 & Lebar / Tinggi & \multicolumn{2}{c}{0,30 m / 0,60 m (1,08 m$^3$)} \\
  & Tulangan utama & \multicolumn{2}{c}{$2\phi28/2\phi22/2\phi28/3\phi29/2\phi28/2\phi22$ (170,3307 kg)} \\
  & Tulangan sengkang & \multicolumn{2}{c}{$\phi10$--180} \\
\midrule
\multicolumn{2}{l}{\textbf{Total Volume Beton}} & \multicolumn{2}{c}{\textbf{8,635 m$^3$}} \\
\multicolumn{2}{l}{\textbf{Total Berat Besi}} & \multicolumn{2}{c}{\textbf{1275,336 kg}} \\
\end{longtable}
\endgroup

\begin{table}[htbp]
\centering
\caption{Rincian Harga Struktur Portal Ruang}
\label{tab:4-6}
\begin{tabular}{lccc}
\toprule
Material & Harga Satuan (Rp.) & Pemakaian Material & Harga (Rp.) \\
\midrule
Beton & 250.000,00/m$^3$ & 8,635 m$^3$   & 2.158.750 \\
Baja  & 5.000,00/kg      & 1275,336 kg   & 6.376.680 \\
\midrule
\textbf{Total} & & & \textbf{8.535.430} \\
\bottomrule
\end{tabular}
\end{table}
